\section{Limits}

\subsection{First examples of (co)limits}

Our first example of a limit is a {\bf product}.
Suppose we have objects $X,Y\in\C$, their product is an object $X\times Y$ with associated morphisms $p_X\colon X\times Y\to X$ and $p_Y\colon X\times Y\to Y$ such that
for any other object $C\in\C$ with $q_X\colon C\to X$ and $q_Y\colon C\to Y$, there exists a unique $f\colon C\to X\times Y$ making the following diagram commute:

\centerline{\drawdiagram{
&&$X$\cr
$C$&$X\times Y$\cr
&&$Y$
}{
    \diagarrow{from={2,1}, to={2,2}, dashed, text=$f$, y distance=.25cm}
    \diagarrow{from={2,1}, to={1,3}, text=$q_X$, y distance=.25cm}
    \diagarrow{from={2,1}, to={3,3}, text=$q_Y$, y distance=-.25cm}
    \diagarrow{from={2,2}, to={1,3}, text=$p_X$, y distance=.25cm}
    \diagarrow{from={2,2}, to={3,3}, text=$p_Y$, y distance=-.25cm}
}}

Such an object $X\times Y$, if it exists, is unique up to isomorphism.
This isomorphism is in fact the unique isomorphism which commutes (as above) with $p_X,p_Y$.

This construct can clearly be generalized to products of families of objects, $\set{X_i}_{i\in I}$.
Here are some examples:

\benum
    \item For objects $\set{x_i}_{i\in I}$ in a partial order $P$ (a partial order viewed as a category has a unique morphism $x\to y$ iff $x\leq y$), their product is their infimum (if it exists).
    Indeed, it must be an object $p\in P$ such that for every $c\leq x_i$, $c\leq p$.
    \item In $\Set$, the product is the usual set-theoretic product, and the morphisms $p_i\colon\prod_{i\in I}X_i\to X_i$ are the projections.
    \item In $\grp,\ab,\ring,\Vec_k$, etc.\ the product is the usual product of such algebraic structures (i.e.\ the set-theoretic product with the usual operations).
    \item In $\top$, the category of topological spaces, the product of topological spaces $X_i$ is their set-theoretic product endowed with the {\it product topology}.
    This topology has a basis $\prod_iU_i$ for $U_i\subseteq X_i$ open, with all but finitely many being equal to $X_i$.
    \item In the category of metric spaces, the product of $X_i$ is their set-theoretic product with the supremum metric $d((x_i)_i,(y_i)_i)=\sup_{i\in I}d(x_i,y_i)$.
    Of course this only exists if the metric spaces are uniformly bounded, or we allow for infinite distance.
    \item The product in the category of (small or locally small) categories is the usual product of categories.
\eenum

Our first example of a colimit is a {\bf coproduct}.
It is the dual of a product: the coproduct of $X,Y\in\C$ is an object $X\coprod Y\in\C$ with morphisms $p_X\colon X\to X\coprod Y$ and $p_Y\colon Y\to X\coprod Y$ such that
for every other object $C$ with $q_X\colon X\to C$ and $q_Y\colon Y\to C$, there exists a unique $f\colon X\coprod Y\to C$ such that the dual diagram to the one above commutes.
That is, $fp_x=q_X$ and $fp_Y=q_Y$.
Equivalently, the coproduct of objects is their product in $\C^\op$.

\benum
    \item The coproduct of objects in a partial order is their supremum.
    \item The coproduct in $\Set$ is the {\it disjoint union}: $\coprod_{i\in I}X_i$, which is the union of the sets $X_i\times\Set i$.
    \item The coproduct in $\ab$ is the direct sum: $\bigoplus_{i\in I}X_i$ the subset of $\prod_{i\in I}X_i$ for which all but finitely many elements are $0$.
    \item Coproducts in $\grp$ is the free product (not discussed here).
    \item Coproducts in $\top$ is the set-theoretic disjoint union $\coprod_{i\in I}X_i$ whose topology has a base given by the open sets in each $X_i$.
    \item In the category of commutative rings, the coproduct of $R,S$ is their tensor product $R\otimes S$ with multiplication given by $(r\otimes s)\cdot(r'\otimes s')=(rr)\otimes(ss')$.
\eenum

\subsection{In general}

We can generalize the notion of products and coproducts to {\bf limits} and {\bf colimits}.

\bdefn

    For a (small) category $I$, we define its {\emphcolor initial completion} and {\emphcolor terminal completion} as $I_i$ and $I_f$.
    In both $I_i$ and $I_f$ we add a new object to $I$, call it $*$.
    In $I_i$ we add a unique morphism $*\to x$ for every $x\in I$, and for $I_f$ a unique morphism $x\to *$ for every $x\in I$ (this uniquely determines composition).
    Clearly $*$ is an initial object in $I_i$ and final in $I_f$.

\edefn

\bdefn

    Given a (small) diagram (a functor) $F\colon I\to\C$, a {\emphcolor cone} over $F$ is a functor $\hat F\colon I_f\to\C$ which agrees with $F$ on $I$.
    A morphism of cones $\hat F\to\hat F'$ is a natural transformation $\alpha\colon\hat F\To\hat F'$ such that $\alpha_i={\rm id}_{Fi}$ for every $i\in I$.
    This forms a category of cones, $\C_{/F}$.

    Dually a {\emphcolor cocone} under $F$ is a functor $\hat F\colon I_i\to\C$ which agrees with $F$ on $I$, and morphisms are defined similarly.
    This forms a category of cocones, $\C_{F/}$.

\edefn

Thus a cone over $F$ is simply an object $\hat F(*)=X$ and morphisms $\hat F(*\to i)=f_i\colon X\to X_i$ (where $Fi=X_i$) such that for every $\phi\colon i\to j$ in $I$, we have
$F(\phi)s_i=s_j$.
So we can view a cone simply as this data $(X,s_i)_{i\in I}$.
A morphism of cones $(X,s_i)\to(Y,t_i)$ is simply a morphism $f\colon X\to Y$ such that $s_i=t_if$ for every $i\in I$.

Dually a cocone can be viewed as $(X,s_i)_{i\in I}$ such that for every $\phi\colon i\to j$ in $I$, we have $s_jF(\phi)=s_i$.
A morphism of cocones $(X,s_i)\to(Y,t_i)$ is simply a morphism $f\colon X\to Y$ such that $t_i=fs_i$.

\bdefn

    A {\emphcolor limit} of $F$ is a terminal object in $C_{/F}$, and a {\emphcolor colimit} is an initial object in $C_{F/}$.

\edefn

Since terminal and initial objects are unique up to unique isomorphism, the limits of a functor (if they exist) are all isomorphic (as are colimits), with a unique
isomorphism that commutes with the diagram.

We denote the (we can say ``the'' since it is unique up to isomorphism) limit of $F$ by $\lim_IF$ and its colimit by $\colim_IF$.
Formally these are objects of $\C$ along with extra information, but generally we keep that information implicit, and write $\lim_IF,\colim_IF\in\C$ to be their associated objects.

We note that a colimit is just a limit in the opposite category.
Explicitly, a cocone under $F\colon I\to\C$ is the same as a cone over $F^\op\colon I^\op\to\C^\op$: $I^\op_f=(I_i)^\op$ and so a cone over $F^\op$ is a functor
$\hat F^\op\colon(I_i)^\op\to\C^\op$ which agrees with $F^\op$ on $I$, equivalently a functor $\hat F\colon I_I\to\C$ which agrees with $F$ on $I$.
That is to say $\C^\op_{/F^\op}\cong\C_{F/}$.

\subsection{More examples}

Let us consider the index category $I=\bul\tto\bul$ (two objects, two non-identity morphisms), and diagrams over it.
Such a diagram is $X\xvarrightarrows{f}[g]Y$.
Explicitly a cone over this diagram is an object $C$ and $h\colon C\to X,h'\colon C\to Y$ such that $fh=gh=h'$.
Thus $h'$ is redundant and we can view a cone as simple $(C,h)$ where $h\colon C\to X$ such that $fh=gh$.

The limit of this diagram is called an {\bf equalizer}: it is $(E,h)$ with $h\colon E\to X$ such that $fh=gh$ and for every other cone $(C,h')$
there is a unique morphism $\phi\colon C\to E$ such that $h'=h\phi$.

For example, in $\Set$ the equalizer is just $\set{x\in X}[fx=gx]\subseteq X$.
This is clearly a cone, and if $fh=gh$ for any other $h\colon C\to X$ we have that $h$ restricts to being in the equalizer, so $\phi$ is just the inclusion.

\subsection{Limits as products and equalizers}

Every limit (in any category) can be written using products and equalizers.
Indeed, let $F\colon I\to\C$ be a diagram and let us consider the following diagram:
$$ f,g\colon\prod_{i\in I}F(i)\tto\prod_{\alpha\colon i\to j\in I}F(j) $$
we construct $f,g$ and show that the limit of $F$ is the equalizer of the above diagram.

We begin by defining $f$.
In general $\hom(\bul,\prod_{i\in I}X_i)\cong\prod_{i\in I}\hom(\bul,X_i)$ naturally, since $f\colon Y\to\prod_{i\in I}X_i$ corresponds to $f_i=p_if$, and $f_i$ gives us a unique $f$ which
factors through the $f_i$s by the universal property of products.
In particular, to define $f$ we need to define $f_\alpha\colon\prod_{i\in I}F(i)\to F(j)$ for every $\alpha\colon i\to j$ to simply be $p_j$.
And similarly we define $g_\alpha\colon\prod_{i\in I}F(i)\to F(j)$ to be the composition $F(\alpha)p_i$.

\bthrm

    If $\phi\colon X\to\prod_{i\in I}F(i)$ is an equalizer for the above diagram, then $(\phi_i\colon X\to F(i))$ is a limit to the diagram $F$.

\ethrm

\bproof

    Since $\phi$ is an equalizer for $f,g$, i.e.\ $f\phi=g\phi$ then for $\alpha\colon i\to j$ we compose with $p_j$ to get $p_j\phi=F(\alpha)p_i\phi$, i.e.\ $\phi_j=F(\alpha)\phi_i$.
    This is exactly what it means for $(\phi_i)$ to be a cone for $F$ though.
    So $F$-cones can be viewed as equalizer cones, and then equalizer morphisms are the same as $F$-cone morphisms.
    Indeed, if $\psi\colon Y\to\prod_{i\in I}F(i)$ is another equalizer cone, then a morphism $\psi\to\phi$ is a morphism $f\colon Y\to X$ such that $\psi=\phi f$.
    That is equivalent to $\psi_i=\phi_if$ for all $i$, i.e.\ a morphism of $F$-cones.
    So we see that equalizers and $F$-limits are the same, as desired.
    \qed

\eproof

Note that the limit of $F$ can exist without $\prod_{i\in I}F(i)$ or $\prod_{\alpha\colon i\to j}F(j)$ existing.
But if a category has arbitrary products and equalizers, then it has all limits.

Similarly for cocones, by dualizing, we get that a coequalizer (since $I^\op=I$) of $X\tto{f}[g]Y$ is an $h\colon Y\to E$ such that $hf=hg$, and is universal in this regard.
And so given a diagram $F\colon I\to\C$, a limit is a coequalizer of
$$ f,g\colon\coprod_{\alpha\colon i\to j}F(j)\to\coprod_{i\in I}F(i) $$
where $f_\alpha\colon F(j)\to\coprod_{i\in I}F(i)$ is the inclusion $\iota_j$ and $g_\alpha\colon F(j)\to\coprod_{i\in I}F(i)$ is the composition $F(\alpha)\iota_i$.

\bdefn

    A category $\C$ is {\emphcolor complete} if it has all (small) limits.
    A category $\C$ is {\emphcolor cocomplete} if it has all (small) colimits.

\edefn

\bcoro

    A category is complete iff it has arbitrary products and equalizers.
    A category is cocomplete iff it has arbitrary coproducts and coequalizers.

\ecoro

\bexam

    In $\Set$, we already have a characterization of products, coproducts, equalizers, and coequalizers.
    Thus the limit of $F\colon I\to\Set$ is simply (since for $\alpha\colon i\to j$, $f_\alpha(x_i)_{i\in I}=x_j$ and $g_\alpha(x_i)_{i\in I}=F(\alpha)(x_i)$),
    $$ \lim_IF = \set{(x_i)_{i\in I}\in\prod_{i\in I}F(i)}[\forall\alpha\colon i\to j.\,x_j=F(\alpha)(x_i)] $$
    The coequalizer of $f,g$ in $\Set$ is $Y/{\sim}$ where $\sim$ is the smalleste equivalence relation such that $f(x)\sim g(x)$ for all $x\in X$.
    Thus
    $$ \colim_IF = \coprod_{i\in I}F(i)/(\forall\alpha\colon i\to j.\,x_j\sim F(\alpha)x_i) $$

\eexam

Since we already know that the following categories have products and coproducts, and we can find equalizers and coequalizers:

\bcoro

    $\Set,\ab,\top,\grp,\ring,{\sf CRing},\Vec_k$ are all complete and cocomplete.

\ecoro

Note that the restriction of (co)complete categories to {\it small\/} limits and colimits is necessary:

\bprop

    If $\C$ is a category such that every diagram (not just small ones) has a limit, then $\C$ is a preorder (a partial order which is not necessarily antisymmetric).

\eprop

\bproof

    We need to show that between every two objects there is at most a single morphism.
    Suppose that there are two morphisms $f,g\colon X\tto Y$, then let $\kappa$ be the cardinality of all the morphisms in $\C$.
    Consider then $Y^\kappa$ (the product $\prod_\kappa Y$), then since $\hom(X,Y^\kappa)\cong\hom(X,Y)^\kappa$, we have that there are at least $2^\kappa$ morphisms $X\to Y^\kappa$,
    a contradiction since $2^\kappa>\kappa$.
    \qed

\eproof

\bprop

    If $I$ has an initial element $i_0$, then for every $F\colon I\to\C$, $\lim_IF=F(i_0)$.
    Dually if $I$ has a terminal element, then $\colim_IF=F(i_0)$.

\eprop

\bproof

    A cone over $F$ is $(X,s_i\colon X\to X_i)$ such that $F(\phi)s_i=s_j$ for all $\phi\colon i\to j$.
    Let $\phi\colon i_0\to i$ be the unique morphism, then we have $s_i=F(\phi)s_{i_0}$ and so we can view a cone simply as $(X,s_{i_0}\colon X\to X_{i_0})$.
    Thus $X_{i_0}$ is a cone ($s_{i_0}=\id$) and it clearly is terminal, thus it is a limit.
    \qed

\eproof

\bexam

    A {\emphcolor directed set} is a partial order $(I,\leq)$ such that every finite set has an upper bound; equivalently for every $a,b\in I$ there exists a $c\in I$ such that $a,b\leq c$.
    The colimit of a functor $I\to\C$ is called a {\emphcolor direct limit}, while the limit of a functor $I^\op\to\C$ is called a {\emphcolor inverse limit}.

    A functor $I\to\C$ is also called a {\emphcolor direct system}, it can be characterized as a collection of objects $\set{X_i}_{i\in I}$ and morphisms $\set{f_{ij}\colon X_i\to X_j}_{i\leq j}$ such that
    $f_{ii}=\id_{X_i}$ and $f_{ik}=f_{jk}\circ f_{ij}$.
    In algebraic categories like $\grp,\ring$, etc.\ we can characterize the direct limit as
    $$ \colim X_i = \bigsqcup_{i\in X_i}X_i/\set{x_i\sim x_j\iff\exists i,j\leq k.\,f_{ik}(x_i)=f_{jk}(x_j)} $$
    where $\sqcup$ is the set-theoretic disjoint union.

    A functor $I^\op\to\C$ is also called a {\emphcolor inverse system}, a collection of objects $\set{X_i}_{i\in I}$ and morphisms $\set{f_{ij}\colon X_j\to X_i}_{i\leq j}$ such that
    $f_{ii}=\id_{X_i}$ and $f_{ik}=f_{ij}\circ f_{jk}$ for $i\leq j\leq k$.
    In algebraic categories the inverse limit is
    $$ \lim X_i = \set{(x_i)_{i\in I}\in\prod_{i\in I}X_i}[x_i=f_{ij}(x_j)\hbox{ for all $i\leq j$}] $$

    In particular if $f_{ij}\colon X_i\inj X_j$ are inclusions then $\colim X_i$ is simply the union of the $X_i$s.
    And if $f_{ij}\colon X_j\inj X_i$ are inclusions then $\lim X_i$ is the intersection.

    We consider now the directed set $\omega$ (finite numbers), a functor $\omega\to\C$ is characterized simply by morphisms $X_i\to X_{i+1}$.
    More interestingly consider the maps $\bZ\xto{\times p}\bZ$ for prime $p$.
    The direct limit of this direct system ($\bZ\xto{\times p}\bZ\xto{\times p}\cdots$) is $\bZ[1/p]\subseteq\bQ$, the subgroup of all rational numbers whose denominator is a power of $p$.
    The projections $\bZ_i\to\bZ[1/p]$ (where $\bZ_i$ is the $i$th copy of $\bZ$ in the system) are given by $\times1/p^i$.

    An inverse system of $\omega$ are morphisms $X_{i+1}\to X_i$.
    Let us consider the inverse system whose morphisms $\bZ/p^{k+1}\bZ\to\bZ/p^k\bZ$ are inclusion.
    The inverse limit of this inverse system is called $\bZ_p$, the {\emphcolor $p$-adic integers}, which are $(x_n)_n$ such that $x_n\in\bZ/p^n\bZ$ and $x_n\equiv x_{n+1}\pmod{p^n}$.

    Another example of an interesting inverse limit is the inverse limit of the $\omega^\op$-system of morphisms $S^1\to S^1$ given by wrapping the circle $p$ times around itself.
    This gives an interesting compact Hausdorff topological space.

\eexam

\subsection{Group actions}

A group $G$ can be viewed as a category $\B G$ with a single object whose hom-set is $G$.
Then a $G$-set is just a functor $F\colon\B G\to\Set$ (i.e.\ the two categories are isomorphic).
Indeed, a functor $F\colon\B G\to\Set$ is a set $X$ along with automorphisms $f_g$ for every $g\in G$ which preserve composition.

Since $\B G$ is a small category, we can consider the limits and colimits of $G$-actions.
These are the invariants and coinvariants of the action, respectively:
$$ \lim F = \set{x\in X}[\forall g\in G.\,gx=x] = X^G,\qquad \colim F = X/\set{\forall g\in G.\,gx\sim x} = X_G $$

If we swap $\Set$ with $\top$, then a functor $F\colon\B G\to\top$ is just a topological space $X$ and homeomorphisms for every $g\in G$ which preserve composition.
The limit of such a functor is still the subspace of $G$-invariants, while the colimit is the space of paths with the quotient topology $X\to X/G$.

And if we swap $\Set$ with $\Vec_k$, then a functor $F\colon\B G\to\Vec_k$ is a $G$-representation.
The limit is still the subspace of $G$-invariants, and the colimit is $X/\lspan\set{x-gx}[g\in G,x\in X]$.

\subsection{Pullbacks and pushforwards}

\bdefn

    A {\emphcolor pullback} is the limit of a functor from the category $\bul\to\bul\ot\bul$.
    A {\emphcolor pushforward} is the colimit of a functor from the category $\bul\ot\bul\to\bul$ (the opposite category of that for pullbacks).

\edefn

Since the category $\bul\to\bul\ot\bul$ has a terminal object, colimits over it are noninteresting, hence why we don't define ``copullbacks''.
Similarly $\bul\ot\bul\to\bul$ has an initial object, so we don't define ``copushforwards''.

The limit cone of a pullback is also called a {\bf pullback square}.
That is, a pullback square is

\centerline{\drawdiagram{
$P$&$Y$\cr
$X$&$Z$
}{
    \diagarrow{from={1,1}, to={1,2}}
    \diagarrow{from={1,1}, to={2,1}}
    \diagarrow{from={2,1}, to={2,2}, text=$f$, y distance=.25cm}
    \diagarrow{from={1,2}, to={2,2}, text=$g$, x distance=.25cm}
}}

\noindent where $P$ is the pullback of the diagram.
In $\Set$, the pullback of $X\xto{\ f\ }Z\xot{\ g\ }Y$ is $X\times_ZY=\set{(x,y)\in X\times Y}[f(x)=g(y)]$.

Similarly the colimit cone of a pushforward is called a {\bf pushforward square}:

\centerline{\drawdiagram{
$Z$&$Y$\cr
$X$&$P$
}{
    \diagarrow{from={1,1}, to={1,2}, text=$g$, y distance=.25cm}
    \diagarrow{from={1,1}, to={2,1}, text=$f$, x distance=.25cm}
    \diagarrow{from={2,1}, to={2,2}}
    \diagarrow{from={1,2}, to={2,2}}
}}

\noindent In $\Set$ the pushforward of $X\xot{\ f\ }Z\xto{\ g\ }Y$ is $X\cup_ZY=X\dcup Y/(f(z)\sim g(z))$.

\bprop[title=Gluing lemma]

    Given the following commutative diagram:

    \centerline{\drawdiagram{
        $A$&$B$&$C$\cr
        $A'$&$B'$&$C'$
    }{
        \diagarrow{from={1,1}, to={1,2}, text=$p$, y distance=.25cm}
        \diagarrow{from={1,2}, to={1,3}, text=$q$, y distance=.25cm}
        \diagarrow{from={2,1}, to={2,2}, text=$p'$, y distance=.25cm}
        \diagarrow{from={2,2}, to={2,3}, text=$q'$, y distance=.25cm}
        \diagarrow{from={1,1}, to={2,1}, text=$f$, x distance=.25cm}
        \diagarrow{from={1,2}, to={2,2}, text=$g$, x distance=.25cm}
        \diagarrow{from={1,3}, to={2,3}, text=$h$, x distance=.25cm}
    }}

    \benum
        \item If the right square is a pullback square, then the left square is a pullback square iff the outer square is a pullback square.
        \item If the left square is a pushforward square, then the right square is a pushforward square iff the outer square is a pushforward square.
    \eenum

\eprop

\blemm

    If $\C$ has a terminal object $1$, then the pullback of $A\to1\ot B$ defines the product $A\times B$.
    Dually if $\C$ has an initial object $0$, then the pushforward of $A\ot0\to B$ defines the coproduct $A\sqcup B$.

\elemm

\bproof

    A cone over $A\to1\ot B$ is the data $(X,p_1\colon X\to A,p_2\colon X\to B,!\colon X\to1)$ which commute.
    But since $1$ is terminal, this is equivalent to $(X,p_1\colon X\to A,p_2\colon X\to B)$, i.e.\ a cone over $A,B$.
    Thus pullbacks of this diagram are simply products.
    \qed

\eproof

\blemm

    If $\C$ has binary products then the pullback of $A\xto{(f,g)}B\times B\xot{(1_B,1_B)}B$ is the equalizer of $f,g\colon A\tto B$.
    Dually for pushforwards and binary coproducts.

\elemm

For $f\colon X\to A$ and $g\colon X\to B$ we denote by $(f,g)$ the unique morphism $(f,g)\colon X\to A\times B$ which projects onto $f$ and $g$.

\bproof

    A cone over $A\xto{(f,g)}B\times B\xot{(1_B,1_B)}B$ is the data $(X,\phi\colon X\to A,\psi\colon X\to B,\chi\colon X\to B\times B)$ such that
    $$ (f,g)\circ\phi = (1_B,1_B)\circ\psi = \chi $$
    so equivalently $(X,\phi,\psi)$ such that $(f,g)\phi=(1_B,1_B)\psi$, which is equivalent to requiring that $f\phi=\psi,g\phi=\psi$, so equivalently a cone
    is just $(X,\phi\colon X\to A)$ such that $f\phi=g\phi$, which is a cone over the equalizer diagram of $f,g$.
    \qed

\eproof

By considering our proof from before relating products and equalizers to all (small) limits, we see that we can write all finite limits with finite products and equalizers.
Similar for colimits.
Further note that the finite product $X_1\times\cdots\times X_n$ is equal to $X_1\times(X_2\times\cdots)$ (if all binary products exist).
In general if $\C$ has all products up to a certain size, then for a set $I$ and a partition of $I$, $\set{I_j}_{j\in J}$ we have that $\prod_{j\in J}\prod_{i\in I_j}X_i$
is the product $\prod_{i\in I}X_i$.
We will omit the proof of that here, but it means that a category has all finite products iff it has all binary products.
Similar for coproducts.

Furthermore, consider the empty category $0$.
A functor $F\colon0\to\C$ has a limit iff $\C$ has a terminal object.
Indeed, $\C_{/F}$ is just $\C$, so a limit of $F$ (a terminal cone) is just a terminal object in $\C$.
Since $0$ is a discrete category (i.e.\ no non-trivial morphisms), we can view a terminal object as the empty product.
Similarly an initial object is the empty coproduct.

\bcoro

    A category has all finite limits iff it has all pullbacks and a terminal object.
    Dually a category has all finite colimits iff it has all pushforwards and an initial object.

\ecoro

\subsection{Preservation of (co)limits}

\bdefn

    Let $I$ be an index category, and $c\in\C$ an object.
    We define the {\emphcolor constant functor} at $c$ to be the functor $c\colon I\to\C$ which maps all objects in $I$ to $c$ and all morphisms to $1_c$.
    This defines the so-called {\emphcolor diagonal functor} $\Delta\colon\C\to\C^I$ which maps $c\in\C$ to the constant functor at $c$, and $f\colon c\to d$ to be the natural
    transformation whose components at each $i\in I$ to be $f$.

\edefn

Recall that a cone over $F\colon I\to\C$ is a functor $\hat F\colon I_f\to\C$ which agrees with $F$.
That is, it is the data $(c,f_i\colon c\to F(i))_{i\in I}$ such that for every $\phi\colon i\to j\in I$ we have that $F(\phi)f_i=f_j$.
Notice that this is equivalent to a natural transformation $\lambda\colon c\To F$ (from the constant functor to $F$): $\lambda_i=f_i$.
So we can view cones over $F$ to simply be natural transformations $\lambda\colon c\To F$.
A morphism of cones from $\lambda\colon c\To F$ to $\mu\colon d\To F$ is a $\C$-morphism $f\colon c\to d$ such that $\mu f=\lambda$.
Dually we can consider cocones under $F$ to be natural transformations $\lambda\colon F\To c$.

Consider a diagram $F\colon I\to\C$ and a functor $G\colon\C\to\D$.
This induces a diagram $G\circ F\colon I\to\D$, and as such we can try and relate $\lim_IG\circ F$ and $G(\lim_IF)$.

First notice that $G\colon\C\to\D$ induces a functor $G\colon\C_{/F}\to\D_{/F}$ by whiskering: $\lambda\colon c\To F$ maps to $G\lambda\colon Gc\To G\circ F$.
Similarly it induces a functor $G\colon\C_{F/}\to\D_{F/}$ by whiskering as well: $\lambda\colon F\To c$ maps to $G\lambda\colon G\circ F\To Gc$.
This leads us to the following definition:

\bdefn

    Let $K\colon I\to\C$ be a class of diagram (i.e.\ a collection of functors $I\to\C$) and $G\colon\C\to\D$ be a functor.
    Then
    \benum
        \item $G$ {\emphcolor preserves} $K$'s limits if $\lambda$ is a limit cone of $K$ (of some functor in $K$), $G\lambda$ is a limit cone of $G\circ K$.
        \item $G$ {\emphcolor reflects} $K$'s limits if for every for every cone $\mu\in\C_{/K}$, if $G\mu$ is a limit cone of $G\circ K$, then $\mu$ is a limit cone of $K$.
        \item $G$ {\emphcolor creates} $K$'s limits if it reflects limits and if $G\circ K$ has a limit then there is a limit $\lambda$ of $K$ such that $G\lambda$ is a limit.
    \eenum
    This dualizes to colimits.

\edefn

So if $F\colon\C\to\D$ creates limits for some class of diagrams and $\D$ has limits of those diagrams, then $\C$ has all those limits and $F$ preserves them.

As discussed, $G\colon\C\to\D$ induces a functor $G\colon\C_{/F}\to\D_{/F}$, so if $\lim_IG\circ F$ exists then it is terminal in $\D_{/F}$ and since $G(\lim_IF)\in\D_{/F}$ we have
a unique morphism $\alpha\colon G(\lim_IF)\to\lim_I(G\circ F)$ (it is unique in $\D_{/F}$).
Similarly we have a unique morphism $\beta\colon\colim_I(G\circ F)\to G(\colim_IF)$.
These are called the {\emphcolor assembly maps}.

The following is immediate:

\bprop

    $G$ preserves $F$'s limit (colimit) iff $\alpha$ (respectively $\beta$) is an isomorphism.

\eprop

\bexam

    Consider the {\emphcolor forgetful functor} $U\colon\ab\to\Set$ which maps Abelian groups to their underlying sets and morphisms to their underlying pure functions.
    Given two Abelian groups $A,B$ we have projections $A\xot{p_A}A\times B\xto{p_B}B$.
    Applying $U$ gives us the diagram $U(A)\xot{U(p_A)}U(A\times B)\xto{U(p_B)}U(B)$, which is precisely a cone over $\set{U(A),U(B)}$.
    Thus we have a unique morphism (unique as a cone morphism):
    $$ U(A\times B) \to U(A)\times U(B) $$
    This is an isomorphism, and as such $U$ preserves binary products.

    Similarly we have the coproduct $A\xto{j_A}A\oplus B\xot{j_B}B$, which gives us $U(A)\xto{j_A}U(A\oplus B)\xot{j_B}U(B)$ and thus we have a morphism
    $$ U(A)\sqcup U(B) \to U(A\oplus B) $$
    which is not an isomorphism.
    Thus $U$ does not preserve coproducts, and so $U$ doesn't preserve colimits ($U$ may preserve specific colimits though).

\eexam

The following are immediate:

\bprop

    Let $\C,\D$ be complete categories, then $F\colon\C\to\D$ preserves all limits iff it preserves all products and equalizers.
    Dually for colimits.

\eprop

\blemm

    If $G\colon\C\to\D$ and $H\colon\D\to\E$ are functors which preserve all $I$-limits (limits whose diagram is of shape $I$), then $H\circ G$ preserves all $I$-limits as well.
    Similar for colimits.

\elemm

Notice that we can consider the composition for $F\colon I\to\C$:
$$ H\circ G(\lim_IF) \to H(\lim_IG\circ F) \to \lim_I(H\circ G\circ F) $$
This is the assembly morphism for $H\circ G$.
As such if $H\circ G$ preserves limits and $H$ reflects isomorphisms (i.e.\ $Hf$ is an isomorphism iff $f$ is), then $G$ preserves limits.

\bexam

    If we consider the forgetful functors $\ab\xto{U_1}\grp\xto{U_2}\Set$, then since $U_2$ reflects isomorphisms and $U_2\circ U_1$ preserves limits, $U_1$ preserves limits as well.

\eexam

\bprop

    Fully faithful functors reflect limits and colimits.

\eprop

\bproof

    Let $F\colon I\to\C$ be a diagram and $G\colon\C\to\D$ be fully faithful.
    We claim that the induced functor $G\colon\C_{/F}\to\D_{/G\circ F}$ is fully faithful as well.
    Indeed, let $\lambda,\mu\in\C_{/F}$ be cones of objects $c,d$ respectively; then $f\in\hom(G\lambda,G\mu)$ is a morphism $Gc\to Gd$ such that $(G\mu)f=G\lambda$.
    Clearly $G\colon\hom(\lambda,\mu)\to\hom(G\lambda,G\mu)$ is injective (since it is a restriction of $\hom(c,d)\to\hom(Gc,Gd)$).
    And since $G$ is full, there is a $g\colon c\to d$ such that $Gg=f$; we claim that $g$ is a morphism $\lambda\to\mu$.
    Indeed: we have $(G\mu)(Gg)=G(\mu g)=G\lambda$, and since $G$ is fully faithful this means $\mu g=\lambda$.

    This means that if $G\lambda$ is a limit cone, then for every $\mu\in\C_{/F}$ we have $\hom(\mu,\lambda)\cong\hom(G\mu,G\lambda)=\Set*$.
    Thus $\lambda$ is terminal and thus a limit cone.
    So $G$ reflects limits.
    \qed

\eproof

\bcoro

    If $F\colon\C\to\D$ is an equivalence, it preserves and creates all limits and colimits.

\ecoro

\bproof

    We claim that for any diagram $D\colon I\to\C$, $F\colon\C_{/D}\to\D_{/F\circ D}$ is an equivalence.
    We just need to show essential surjectivity, let $\lambda\colon d\To F\circ D$ be a cone.
    Since $F$ is an equivalence, $\iota\colon d\cong Fc$ for some $c\in\C$ and for each $i\in I$ there exists $\mu_i\colon c\to D(i)$ such that $\lambda_i=F(\mu_i)\iota$.
    Then $\mu$ is a cone $c\To D$ such that $(F\mu)\iota=\lambda$, i.e.\ $F\mu$ and $\lambda$ are isomorphic.

    Since $F$ is fully faithful we already know it reflects limits.
    Now if $\lambda$ is a limit cone over $F\circ D$, then it is isomorphic to $F\mu$ so $\mu$ is a limit cone over $D$.
    Thus $F$ creates all limits.

    And if $\lambda$ is a limit cone over $D$, then for every limit cone $\mu$ over $F\circ D$ it is isomorphic to $F\mu'$ so, $\hom(\mu,F\lambda)\cong\hom(F\mu',F\lambda)\cong\hom(\mu',\lambda)$
    a singleton.
    Thus $F\lambda$ is a limit cone over $F\circ D$, so $F$ preserves limits.
    \qed

\eproof

\bcoro

    Equivalent categories have the same shape limits.

\ecoro

\bexam

    We claim that $\Vec_k$ and $\Vec_k^\op$ are not equivalent.
    Note that all vector spaces are isomorphic to coproducts of $k$.
    Now suppose that $E\colon\Vec_k\to\Vec_k^\op$ is an equivalence, then since it preserves coproducts it must map $\coprod_{i\in I}k$ to $\coprod_{i\in I}Ek$.
    But notice that coproducts in the opposite category are simply products, so $\coprod_{i\in I}k$ maps to $\prod_{i\in I}Ek$.
    Now notice that $k\cong\hom(k,k)\cong\hom(Ek,Ek)$, since in general $V$ can be embedded into $\hom(V,V)$ we have that $Ek$ has cardinality at most $\abs k$.

    Take $I$, $J$ index sets, then we have that
    $$ \hom(\coprod_Ik,\coprod_Jk) \cong \prod_I\hom(k,\coprod_Jk) \cong \prod_I\coprod_Jk $$
    and
    $$ \hom(\prod_JEk,\prod_IEk) \cong \prod_I\hom(\prod_JEk,Ek) $$
    in general we have that $\hom(V,Ek)\supseteq V$ (comparing cardinalities), and so we get that the above is a superset of $\prod_I\prod_JEk\supset\prod_{I,J}k$.
    Thus by equivalence we must have that $\abs{\prod_{I,J}k}\leq\abs{\prod_I\coprod_Jk}$.
    But for $I,J$ large enough this means $\abs{k^{I\times J}}\leq\abs J^{\abs I}$.
    Letting $J=k^I$ we get $k^J=k^{k^I}\leq k^I$ which is impossible.

\eexam

\bdefn

    An object $c\in\C$ can be viewed as a functor from the one-object category $c\colon1\to\C$ (there is an isomorphism $\C\cong\C^1$).
    The category of cones over $c$ is called the {\emphcolor overcategory} of $c$ and is denoted $\C/c$ (instead of $\C_{/c}$).
    Dually the category of cocones under $c$ is called the {\emphcolor undercategory} of $c$ and is denoted $c/\C$ (instead of $\C_{c/}$).

\edefn

Now, a cone over $c$ is just a natural transformation $\lambda\colon X\To c$ for some $X\in\C$.
Such a natural transformation is just a morphism $f\colon X\to c$.
A morphism of cones from $f\colon X\to c$ to $g\colon Y\to c$ is a morphism $h\colon X\to Y$ such that $f=gh$.
So we have an alternative characterization of $\C/c$:
\benum
    \item Objects are pairs $(X,f)$ where $X\in\C$ and $f\colon X\to c$ is a morphism;
    \item A morphism $(X,f)\to(Y,g)$ is a morphism $h\colon X\to Y$ such that $f=gh$.
\eenum

Similarly a cocone under $c$ is a morphism $f\colon c\to X$.
A morphism of cocones $f\colon c\to X$ to $g\colon c\to Y$ is a morphism $h\colon X\to Y$ such that $hf=g$.

In general we have a forgetful functor $U\colon\C_{/F}\to\C$ which maps $\lambda\colon c\To F$ to $c$.
And we also have $U\colon\C_{F/}\to\C$ which maps $\lambda\colon F\To c$ to $c$.
Special cases of these are of course forgetful functors from the overcategory and undercategory.

We can characterize the category of cones another way:

\bdefn

    Let $F\colon\C\to\Set$ be a $\Set$-valued functor.
    Then $F$'s {\emphcolor category of elements} of $\C$ is the category $\int F$ whose objects are pairs $(c,x)$ where $c\in\C$ and $x\in F(c)$, and a morphism
    $(c,x)\to(d,y)$ is a $\C$-morphism $f\colon c\to d$ such that $y=F(f)(x)$.

    For a contravariant functor $F\colon\C^\op\to\Set$, it has a category of elements of $\C^\op$.
    We define its category of elements of $\C$ to be the opposite category of this: $\int_\C F=(\int_{\C^\op}F)^\op$.

\edefn

Explicitly, for a contravariant functor $F\colon\C^\op\to\Set$ its category of elements $\int_\C F$ has pairs $(c,x)$ for $c\in\C$ and $x\in F(c)$ and a morphism
$(c,x)\to(d,y)$ is a $\C$-morphism $f\colon c\to d$ such that $x=F(f)(y)$.

Let $F\colon I\to\C$ be a diagram, we define two functors:
$$ \cone(\bul,F)\colon\C^\op\to\Set,\qquad \cone(F,\bul)\colon\C\to\Set $$
where $\cone(c,F)$ is equal to the set of natural transformations $c\To F$, and $\cone(F,c)$ is the set of natural transformations $F\To c$.
For a morphism $f\colon c\to d$, $\cone(f,F)\colon\cone(d,F)\to\cone(c,F)$ by $\lambda\mapsto\lambda\circ f$.
Similarly $\cone(F,f)\colon\cone(F,c)\to\cone(F,d)$ by $\lambda\mapsto f\circ\lambda$.

Now, we can characterize the category of cones and cocones as follows:
$$ \C_{/F} = \int\cone(\bul,F),\qquad \C_{F/} = \int\cone(F,\bul) $$
Indeed, an object in the category of elements of $\cone(\bul,F)$ is a pair $(c,\lambda\colon c\To F)$ and morphisms $(c,\lambda)\to(d,\mu)$ are morphisms $f\colon c\to d$ such that
$\lambda=\cone(f,F)(\mu)=\mu f$.
This is precisely the definition of the category of cones.
Similarly for cocones.

There is a clear forgetful functor $U\colon\int F\to\C$ which maps $(c,x)\mapsto c$.
Note that because of our definition for contravariant functors, this is a forgetful functor both when $F$ is covariant and when it is contravariant.

Let us now consider limits of the category of elements.
To do so let us consider what it means for $D\colon I\to\int F$ to be a diagram.
We can view such diagrams equivalently as diagrams: $UD\colon I\to\C$ as well as a cone $\lambda\colon*\To FUD$ (viewing both as functors $I\to\Set$, where $*$ is a singleton).
Indeed: if $D(i)=(c_i,x_i\in F(c_i))$ we can define $\lambda$ to be $\lambda_i=x_i\in FUD(i)$.
And conversely given a diagram $D\colon I\to\C$ and a cone $\lambda\colon*\to FD$ then this defines a diagram $I\to\int F$ by $i\mapsto(D(i),\lambda_i)$ and $\phi\colon i\to j$ maps
to $D(\phi)$.

So a diagram $I\to\int F$ is simply a diagram $D\colon I\to\C$ as well as a cone $a\colon*\To FD$ in $\Set$.
A morphism from a cone $\mu\colon(d,y)\To(D,a)$ to $\lambda\colon(c,x)\To(D,a)$ is a morphism $f\colon d\to c$ such that $x=Ff(y)$ and $U\mu=(U\lambda)f$.

\bprop

    For a covariant functor $F$, the forgetful functor $U\colon\int F\to\C$ creates all limits which $\C$ admits and $F$ preserves.
    What this means is that for every diagram $D\colon I\to\C$ which has a limit in $\C$ preserved by $F$, then this limit in $\C$ is created by $U$.

\eprop

\bproof

    We need to show that if $D\colon I\to\C$ has a limit preserved by $F$, and $a\colon*\To FD$ is a cone, then:
    \benum
        \item if $\lambda$ is a cone over $(D,a)$ such that $U\lambda$ is a limit of $D$, then $\lambda$ is a limit of $(D,a)$.
        \item there exists a limit $\lambda$ over $(D,a)$ such that $U\lambda$ is a limit of $D$.
    \eenum
    Let $\eta\colon\lim D\To D$ be the limit cone of $D$, then since $F$ preserves it $F\eta\colon F(\lim D)\To FD$ is a limit cone as well.

    Now suppose there is a $\bar\eta\colon X\To(D,a)$ such that $U\bar\eta=\eta$.
    This means that $UX=\lim D$ and $\bar\eta_i=\eta_i$ for all $i\in I$.
    Thus, $\bar\eta\colon(\lim D,x)\To(D,a)$ for $x\in F(\lim D)$, and since it is a cone we must have that $F(\eta_i)(x)=a_i$ for all $i\in I$ as well.
    Now we know that $F(\eta)\colon F(\lim D)\To FD$ is a limit cone, and so there exists a unique morphism from $a\colon*\To FD$ to $F(\eta)$.
    A morphism from $a$ to $F(\eta)$ is precisely a choice of element $x\in F(\lim D)$ such that $F(\eta_i)(x)=a_i$ for all $i\in I$, as desired.

    So there does indeed exist such a $\bar\eta\colon(\lim D,x)\To(D,a)$ such that $U\bar\eta=\eta$.
    We claim that $\bar\eta$ is a limit of $(D,a)$, proving both points.
    Let $\lambda\colon(c,y)\To(D,a)$ be a cone, then $\hom(\lambda,\bar\eta)$ is the set of all morphisms $f\colon c\to\lim D$ such that $x=Ff(y)$ and $\eta=(U\lambda)f$.
    This is a subset of $\hom(U\lambda,\eta)$ which is a singleton.
    Now, the extension of $U\lambda$ to $\lambda$ corresponds to the morphism $a\to F(U\lambda)$ choosing $y$ (as before), and the composition of this with the unique morphism
    $Ff\colon F(U\lambda)\to F(\eta)$ must be equal to the unique morphism $a\to F(\eta)$ corresponding to choosing $x$.
    This means that we have $Ff(y)=x$ as desired.
    So we have that $\bar\eta$ must indeed be terminal, as desired.
    \qed

\eproof

\bcoro

    If $\C$ is complete, then $U\colon\int F\to\C$ creates and preserves all limits $F$ preserves.

\ecoro

If $F$ is contravariant then the forgetful functor $U\colon\int_\C F\to\C$ is simply the opposite functor of $U^\op\colon\int_{\C^\op}F\to\C^\op$ which is the forgetful
functor discussed above.
Since $U^\op$ creates limits, this means that $U$ creates colimits.

\bcoro

    If $F\colon\C^\op\to\Set$ is a contravariant functor then the forgetful functor $U\colon\int_\C F\to\C$ creates all colimits which $\C$ admits and $F$ preserves.

\ecoro

In particular for our cone categories, we must consider what limits and colimits the functors $\cone(F,\bul)\colon\C\to\Set$ and $\cone(\bul,F)\colon\C^\op\to\Set$ preserve.

\bprop

    If $\C$ is complete, then $U\colon\C_{F/}\to\C$ creates and preserves all limits.
    If $\C$ is cocomplete then $U\colon\C_{/F}\to\C$ creates and preserves all colimits.

\eprop

\bproof

    By the above arguments, we just need to show that $\cone(F,\bul)$ and $\cone(\bul,F)$ preserve all limits and colimits respectively.
    Notice that if $\set{A_i}_{i\in I}$ are objects in $\C$ then $\cone(F,\prod_{i\in I}A_i)\cong\prod_{i\in I}\cone(F,A_i)$ naturally, and we also have that
    the cocone of an equalizer is naturally isomorphic to the equalizer of cocones.
    Thus $\cone(F,\bul)$ preserves all limits, and thus by our above arguments $U$ creates and preserves all limits.
    \qed

\eproof

\subsection{Limits in functor categories}

For categories $\C,\D$ we consider the functor category $\D^\C$ of functors $\C\to\D$ whose morphisms are natural transformations.
We alternatively write it by $[\C,\D]$.

\blemm

    For categories $\C,\D,\E$ there is an isomorphism of categories $[\C,[\D,\E]]\cong[\C\times\D,\E]\cong[\D,[\C,\E]]$.

\elemm

\bproof

    We give an isomorphism $\Phi\colon[\C,[\D,\E]]\to[\C\times\D,\E]$.
    For a functor $F\colon\C\to\E^\D$, we define $\Phi(F)\colon\C\times\D\to\E$ by $\Phi(F)(c,d)=F(c)(d)$ on objects and if $\alpha\colon F\To G$ is a natural transformation,
    we note that this means for every $c\in\C$, $\alpha_c\colon Fc\To Gc$ is a natural transformation of functors $\D\to\E$, and so it has a component $\alpha_{c,d}\colon F(c)(d)\to G(c)(d)$.
    So we define $\Phi(\alpha)\colon\Phi(F)\To\Phi(G)$ is given by $\Phi(\alpha)_{(c,d)}=\alpha_{c,d}$.
    This is a functor.

    We define its inverse $\Psi\colon[\C\times\D,\E]\to[\C,[\D,\E]]$ as follows: for $F\colon\C\times\D\to\E$ define $\Psi(F)\colon\C\to\E^\D$ by $\Psi(F)(c)(d)=F(c,d)$ on objects $c\in\C$.
    For morphisms $f\colon c\to c'$ we define $\Psi(F)(f)\colon\Psi(F)(c)\To\Psi(F)(c')$ by $\Psi(F)(f)_d=F(f,1_d)$.
    And for $\alpha\colon F\To G$ natural, we define $\Psi(\alpha)_{c,d}\colon F(c,d)\to G(c,d)$ to be $\alpha_{(c,d)}$.

    The last isomorphism follows because $\C\times\D\cong\D\times\C$.
    \qed

\eproof

This gives us an isomorphism $[\C,[\D,\E]]\cong[\D,[\C,\E]]$ which maps $F\colon\C\to\E^\D$ to $\hat F\colon\D\to\E^\C$ defined by
\benum
    \item $\hat F(d)(c)=F(c)(d)$,
    \item for $f\colon c\to c'$, $\hat F(d)(f)=F(f)_d$,
    \item for $f\colon d\to d'$, $\hat F(f)_c=F(c)(f)$.
\eenum

\bthrm

    Let $F\colon I\to[\C,\D]$ be a diagram and consider its curried functor $\hat F\colon\C\to[I,\D]$.
    Show that if for every $X\in\C$, $\hat F(X)\colon I\to\D$ has a limit $\lim_I\hat F(X)$, then $F$ has a limit and there are natural isomorphisms
    $$ (\lim_I F)(X) \cong \lim_I\hat F(X) $$

\ethrm

\bproof

    We need to show that there is a functor $\lim_IF\colon\C\to\D$ which acts on objects by $(\lim_IF)(X)=\lim_I\hat F(X)$ and that it has an associated terminal cone.
    For brevity we write $L\colon\C\to\D$ instead of $\lim_IF$.

    Let $f\colon X\to Y$ be a $\C$-morphism, then we have a natural transformation $\hat F(f)\colon\hat F(X)\To\hat F(Y)$.
    Now let $\lambda_X\colon\lim_I\hat F(X)\To\hat F(X)$ be the limit cone, then $\hat F(f)\lambda_X\colon\lim_I\hat F(X)\To\hat F(Y)$ is another cone, and as such there is a unique
    $\D$-morphism $L(f)\colon\lim_I\hat F(X)\To\lim_I\hat F(Y)$ such that $\lambda_YL(f)=\hat F(f)\lambda_X$.
    So we define $L$ by:
    $$ L(X) = \lim_I\hat F(X),\qquad \lambda_YL(f)=\hat F(f)\lambda_X $$
    This is functorial: $\hat F(gf)\lambda_X=\hat F(g)\hat F(f)\lambda_X=\hat F(g)\lambda_YL(f)=\lambda_ZL(g)L(f)$, since $L(gf)$ is the unique morphism satisfying this, we have $L(gf)=L(g)L(f)$.

    Now we need to find a cone $\Lambda\colon L\To F$, that is for every $i\in I$ a natural transformation $\Lambda_i\colon L\To F(i)$ such that everything commutes.
    For $X\in\C$ we have the projection $(\lambda_X)_i\colon\lim_I\hat F(X)=L(X)\to\hat F(X)(i)=F(i)(X)$, so define $\Lambda_{i,X}=\lambda_{X,i}$.
    Each $\Lambda_i$ is natural: this assertion just means that $\Lambda_{iY}L(f)=F(i)(f)\Lambda_{iX}$ for all $f\colon X\to Y$.
    Expanding definitions, this just means $\lambda_{Yi}L(f)=\hat F(f)_i\lambda_{Xi}$, which clearly holds by definition of $L(f)$.
    And $\Lambda$ as a whole is natural: to show $\Lambda_j=F(\phi)\Lambda_i$ for every $\phi\colon i\to j$ it is necessary and sufficient to show this for every component, that is to show
    $\lambda_{Xj}=F(\phi)_X\lambda_{Xi}$.
    This is true by naturality of $\lambda_X$ (since $\hat F(X)(\phi)=F(\phi)_X$).

    Now we claim that it is terminal.
    Let $f\colon\C\to\D$ be another functor with a cone $\Phi\colon f\To F$, we need to show that there is a unique natural transformation $\alpha\colon f\To F$ such that $\Phi=\Lambda\alpha$.
    This is equivalent to saying that for every $i\in I,X\in\C$ we have $\Phi_{iX}=\lambda_{Xi}\alpha_X$.
    And indeed this is true: $\Phi_{\bul X}\colon f(X)\To\hat F(X)$ is a cone over $\hat F(X)$ and since $\lambda_X$ is a limit cone there exists a unique $\alpha_X$ such that $\Phi_{\bul X}=\lambda_X\alpha_X$.
    This $\alpha$ is natural and thus concludes our proof.
    \qed

\eproof

This tells us that limits (and by duality colimits) in functor categories can be computed ``pointwise'': if $F\colon I\to[\C,\D]$ is a diagram then we have
$$ (\lim_IF_i)(X) = \lim_I(F_i(X)) $$

\bcoro

    If $\D$ has all (co)limits of shape $I$, then so does the functor category $[\C,\D]$.
    In particular the functor categories over (co)complete categories are all (co)complete.

\ecoro

\subsection{Functorality of limits and colimits}

Notice that if $\C$ has all limits of shape $I$, then we can define a functor
$$ \lim_I\colon[I,\C]\to\C $$
which is defined for functors to be the limit of $F$ (we need choice to choose an explicit representation of the limit in $\C$, as it is only defined up to isomorphism).
For a natural transformation $\alpha\colon F\To G$ of functors $I\to\C$, note that if $\lambda_F\colon\lim_IF\To F$ is the limit cone over $F$, then $\alpha\lambda_F\colon\lim_IF\To G$ is a cone as well,
and as such there is a unique morphism $f_\alpha\colon\lim_IF\to\lim_IG$ such that $\alpha\lambda_F=\lambda_Gf_\alpha$, so define $\lim_I(\alpha)=f_\alpha$.
This is the same construction as before; it is functorial for the same reason.

Similarly if $\C$ has all colimits of shape $I$ we define
$$ \colim_I\colon[I,\C]\to\C $$

\bthrm

    If $F\colon I\times J\to\C$ is a diagram, then considering its curried functors $I\to[J,\C]$ and $J\to[I,\C]$ we have
    $$ \lim_{i\in I}\lim_{j\in J}F(i,j) \cong \lim_{I\times J}F \cong \lim_{j\in J}\lim_{i\in I}F(i,j) $$
    when either the first or last limits exist (and their existence implies the existence of the middle limit).

\ethrm

We omit the proof for now, and will give a slick one in the future.
We also have:

\bthrm

    $\lim_I$ preserves all limits and $\colim_I$ all colimits.

\ethrm

Which we will prove slickly in the future as well.

